\chapter{Introduction}

Nowadays, mathematics is studied on a foundation relying on set theory, that is, all objects of interest are considered as sets. One of the most widely used set theory is ZFC, the Zermelo–Fraenkel set theory (ZF) together with the Axiom of Choice (AC). There are many useful consequences of AC. For example, AC is equivalent to the statement that every vector space has a basis. In set theory, cardinalities of any two sets are comparable if AC is assumed. This is no longer provable when AC is dropped, and many interesting questions on cardinals arise when we study set theory in ZF alone. 

For a set $A$ which is of cardinality $\frak{m}$, the cardinal $\frak{m}!$ is the cardinality of the set of bijections on $A$. Dawson and Howard showed in \cite{DH} that under ZFC, $\frak{m}!=2^{\frak{m}}$ for any infinite cardinal $\frak{m}$. They also showed that each of the statements ``$\frak{m}!<2^{\frak{m}}$ for some infinite cardinal $\frak{m}$", ``$\frak{m}!>2^{\frak{m}}$ for some infinite cardinal $\frak{m}$" and ``$\frak{m}!$ is incomparable with $2^{\frak{m}}$ for some infinite cardinal $\frak{m}$" is relatively consistent with ZF. Roughly speaking, in ZF we cannot conclude any relationship between $2^{\frak{m}}$ and $\frak{m}!$ for an arbitrary infinite cardinal $\frak{m}$.
In \cite{HalSheCon}, Halbeisen and Shelah studied the relationships between $2^{\frak{m}}$ and $\seq(\frak{m})$, the cardinality of the set of finite sequences of elements in $A$ (as well as $\seq^{1-1}(\frak{m})$, the set of finite sequences of elements in $A$ without repetitions). Sonpanow and Vejjajiva also studied the relationships between $\frak{m}!$ and $\seq(\frak{m})$ (as well as $\seq^{1-1}(\frak{m})$) in \cite{SonVej2019}.

In this thesis, we study the cardinalities of $I(A)$, $J(A)$ and $A^{A}$ for a set $A$ where $I(A)$, $J(A)$ and $A^{A}$ are the set of injections, the set of surjections and the set of functions on a set $A$ respectively. We provide some basic knowledge for our works in Chapter II. Our results are divided into two parts: Results in ZF are discussed in Chapter III, and consistency results are discussed in Chapter IV. Finally, we conclude our results in Chapter V.
